\scanpage{392}
\mspage{116}
% archive sheet number: 391
pose
\[
p_!(\deleted{\unsure{d}}\mathrm{ho}_{W,S}(y))=
\deleted{\unsure{\mathrm{ho}\,p_!}}\mathrm{ho}_W(p_!(y')) .
\]
% slanted
\margin{\unsure{opération|expliciter} $p_!$ au niveau de $M(S)$}

Quand $y=p^*(\unsure{\xi|\zeta})$, la flèche d'adjonction \inserted{\unsure{pertinente}} dans $\mathrm{Ho}\deleted{\unsure{t}}_{\deleted{W}\Sigma}$ \inserted{$=D(e)$}, s'obtient par composition
\begingroup
\makeatletter
\tagsleft@true
\makeatother
\[
p_!(y')\longrightarrow p_!(p^*(\xi))\xrightarrow{\mathrm{adj}}\xi\,,
\tag*{$(\alpha)$}
\]
\endgroup
où ici $p_!$ est encore pris dans le sens ordinaire ($M(S)\unsure{\relbar|\longrightarrow}M(e)=M$). Or $S$ ayant un objet initial $\unsure{o|0}$, on a
\[
\begin{aligned}
p_!(y')&\simeq y'(o)\\
p_!(p^*(\xi))&\simeq p^*(\xi)(o)\simeq\xi\,,
\end{aligned}
\]
et la flèche d'adjonction dans $(\alpha)$ est l'identité, la \deleted{flèche} première flèche est la flèche $y'(o)\longrightarrow y(o)$ induite par $y'\relbar y$, elle est dans $\Sigma$, donc $(\alpha)$ est dans $\Sigma$, \deleted{\unsure{a fortiori}} \inserted{donc} c'est un iso \unsure{dans} $\mathrm{Ho}_{\Sigma}\unsure{=|\simeq}D(e)$, \unsure{qed}.

Le même argument marche d'ailleurs si $M$ est seulement une cat. à cofibrations ou à fibrations, mais en se bornant alors au cas où $S$ est un ens ordonné fini.

